\section{Optimal Control Formulation}
\label{ch:optimalcontrol}

The optimal-control extension is the point at which statistical approximation, numerical transcription, and physical modelling become inseparable. A chance constraint is evaluated on a residual produced by a trajectory, but the trajectory itself is generated by a discretized differential equation. Consequently, a small test violation can coexist with a failed nonlinear program, and a reported objective can depend on the quadrature convention rather than on the physical control effort. This chapter states the formulation used by the local implementation, identifies the differences from the published lunar benchmark, and derives an audit identity that limits the interpretation of the local costs.

\subsection{Stochastic optimal-control problem}

Let $y(t)\in\mathbb R^{n_y}$ and $u(t)\in\mathbb R^{n_u}$ denote state and control trajectories, and let $\xi$ denote uncertain parameters. A generic finite-horizon problem is
\begin{equation}
  \begin{aligned}
    \min_{y,u}\quad &J(y,u)=\int_0^{t_f}L(y(t),u(t))\,\mathrm dt+\Phi(y(t_f)),\\
    \text{s.t.}\quad &\dot y(t)=F(y(t),u(t),\xi),\\
    &y(0)\in\mathcal Y_0,\quad u(t)\in\mathcal U,\\
    &\Prob\{g_\ell(y,u,\xi)\leq0\}\geq1-\varepsilon_\ell,
    \quad \ell=1,\ldots,m.
  \end{aligned}
  \label{eq:ocp-general}
\end{equation}
The residual $g_\ell$ may be a terminal event, a path event evaluated at a time $t$, or a vector event reduced to a scalar statistic. The distinction matters: a terminal probability is not a pathwise probability over a continuum of times. In the implementation, path constraints are evaluated at mesh nodes. They are therefore nodewise surrogate constraints unless an additional union or supremum construction is imposed.

For training samples $\{\xi_i\}_{i=1}^{n}$ and a decision vector $z$ containing the nodal states and controls, the residual observations are
\begin{equation}
 r_{\ell i}(z)=g_\ell(y(z),u(z),\xi_i).
\end{equation}
The empirical chance constraint used in Chapters~3--5 is replaced by an integrated-kernel average,
\begin{equation}
 \widehat p_{\ell,h}(z)=\frac{1}{n}\sum_{i=1}^{n}
 \overline K\!\left(\frac{r_{\ell i}(z)}{h}\right),
 \qquad \widehat p_{\ell,h}(z)\leq\varepsilon_\ell,
 \label{eq:ocp-kde}
\end{equation}
where $\overline K$ is the integrated kernel and $h$ is a bandwidth. Equation~\eqref{eq:ocp-kde} is a deterministic NLP constraint after the samples and bandwidth have been fixed. It does not by itself provide a confidence level for the unknown probability. The biased construction discussed in Chapter~4 can provide an upper bound only when its bias and sampling assumptions are satisfied; the local shifted estimator is recorded as an implementation surrogate and is not identified with the Split--Bernstein construction of \citet{keil2020,keil2021}.

\subsection{Direct transcription used in the experiments}

The local fallback solver uses $N$ uniformly spaced nodes $0=t_0<\cdots<t_{N-1}=t_f$ and forward Euler defects. With $h_j=t_{j+1}-t_j$,
\begin{equation}
 y_{j+1}-y_j-h_jF(y_j,u_j,\xi_0)=0,
 \qquad j=0,\ldots,N-2,
 \label{eq:euler-defect}
\end{equation}
where $\xi_0$ is absent from the deterministic dynamics in the lunar proxy. The collocation vector includes all nodal controls, although the last control does not enter the forward Euler defects. Equality constraints enforce the initial state, the defects, and terminal velocity. Terminal position remains a chance-constrained quantity. Inequality constraints enforce control bounds and the KDE approximations of the terminal and path residuals.

The implementation separates four operations. The problem module defines dynamics and residuals; the collocation module maps a vector $z$ to nodal states and controls; the KDE module evaluates integrated-kernel statistics; and the solver module supplies SLSQP with equality and inequality callables. This separation makes it possible to inspect a statistical estimator without silently changing the dynamics. It also exposes a numerical limitation: the same nodal control can contribute to a path residual and to the objective while not contributing to the state defects when it is the final control. That limitation is quantified below rather than hidden in a single objective number.

The continuation sequence is nominal, worst-case, scenario, unbiased KDE, and local shifted Epanechnikov. A stage is warm-started with the last successful vector. Warm starts are consistent with the continuation strategy described by \citet{keil2021warm}, but the local implementation does not use LGR nodes, adaptive mesh refinement, IPOPT, or GPOPS-II. A failed SLSQP stage is retained in the result ledger with its status and constraint margin; it is not promoted to a feasible solution because its held-out violation happens to be below the target.

\subsection{Local lunar landing model}
\label{sec:lunar-ocp}

The local benchmark has altitude $y_1$, vertical velocity $y_2$, and nonnegative thrust $u$:
\begin{equation}
 \dot y_1=y_2,\qquad \dot y_2=-g+u,\qquad g=1.622,\qquad 0\leq u\leq3.
 \label{eq:lunar-dynamics}
\end{equation}
The local horizon is fixed at $t_f=20$, the initial state is $(y_1(0),y_2(0))=(100,0)$, and only terminal velocity is an equality. Terminal altitude is subject to
\begin{equation}
 \Prob\left(\left|y_1(t_f)-\xi_1\right|-\delta>0\right)\leq\varepsilon_a,
 \quad \xi_1\sim\mathcal N(0,0.1^2),\quad \delta=0.25,\quad \varepsilon_a=0.1.
 \label{eq:lunar-terminal-risk}
\end{equation}
The nodewise path event is
\begin{equation}
 \Prob\left(u(t_j)+\xi_2-3>0\right)\leq\varepsilon_b,
 \qquad \varepsilon_b=0.01,
 \label{eq:lunar-path-risk}
\end{equation}
with the locally documented mixture $\xi_2\sim0.6\mathcal N(-0.25,0.10^2)+0.4\mathcal N(0.35,0.15^2)$. The mixture is a synthetic stress test for a bounded, asymmetric disturbance; it is not a claim that the local data reproduce the distribution used in the published study.

The comparison with Keil et al. must be made at the level of formulation. Their lunar study uses $y_1(0)=10$, $y_2(0)=-2$, free final time, free terminal altitude, terminal velocity zero, and $J=\int u\,dt$; it solves the resulting problem with LGR/GPOPS-II/SNOPT and reports costs around $9.09$ for the chance-constrained methods \citep{keil2020,keil2021}. The local problem fixes a different initial altitude, initial velocity, horizon, mesh, and solver. The two objective values therefore cannot be compared numerically. The local experiment is a pipeline and audit experiment, not a replication.

\subsection{A conservation audit of the local objective}

The local transcription provides a stronger diagnostic than a solver status. Summing the velocity defects over the $N-1$ intervals gives
\begin{equation}
  y_{2,N-1}-y_{2,0}=\sum_{j=0}^{N-2}h_j(-g+u_j).
 \label{eq:velocity-telescope}
\end{equation}
With uniform $h=2$, $y_{2,0}=y_{2,N-1}=0$, and $g=1.622$, every feasible trajectory satisfies
\begin{equation}
 \sum_{j=0}^{N-2}h u_j=g t_f=32.44.
 \label{eq:fuel-invariant}
\end{equation}
Thus, in this fixed-horizon model, the physically consistent piecewise-constant fuel integral is invariant over the feasible set. Trajectory shape still matters because terminal altitude imposes the weighted moment
\begin{equation}
 \int_0^{t_f}(t_f-t)u(t)\,\mathrm dt
 =\frac{g t_f^2}{2}-y_1(0)=224.4,
 \label{eq:altitude-moment}
\end{equation}
when terminal altitude is zero. The stochastic terminal event changes the admissible timing of thrust, not the invariant integral implied by terminal velocity.

The code records a trapezoidal integral over all $N$ controls, whereas the Euler defects use only $u_0,\ldots,u_{N-2}$. On a uniform mesh,
\begin{equation}
 J_{\mathrm{trap}}=h\left(\frac{u_0}{2}+\sum_{j=1}^{N-2}u_j+\frac{u_{N-1}}{2}\right)
 =32.44+\frac{h}{2}(u_{N-1}-u_0).
 \label{eq:quadrature-mismatch}
\end{equation}
The nominal record $J=29.8052$ is explained by the endpoint term: the stored solution has approximately $u_0=2.6345$ and $u_{N-1}=0$. The later records at $32.44$ are not evidence that robust control requires more physical fuel. They reveal that the endpoint contribution becomes different when the optimizer is warm-started under another constraint. This audit is a substantive result of the local implementation and is the reason Chapter~8 reports the lunar objective only as a transcription diagnostic.

\begin{table}[tb]
\centering\small
\caption{Formulation differences between the published lunar study and the local audit benchmark.}
\label{tab:lunar-formulation-comparison}
\begin{tabular}{p{0.27\textwidth}p{0.31\textwidth}p{0.31\textwidth}}
\toprule
Feature & Keil et al. & Local implementation\\
\midrule
Initial state & $(10,-2)$ & $(100,0)$\\
Final time & free & fixed at $20$\\
Terminal conditions & velocity equality; altitude free & velocity equality; altitude chance event\\
Transcription & LGR with adaptive machinery & 11-node forward Euler\\
Solver & GPOPS-II/SNOPT & SciPy SLSQP\\
Uncertainty & terminal normal and path mixture in published model & synthetic normal and two-mode mixture\\
Interpretation of $J$ & physical $\int u\,dt$ & trapezoidal nodal diagnostic with endpoint mismatch\\
\bottomrule
\end{tabular}
\end{table}

\subsection{Stage logic and numerical interpretation}

The local stage ledger contains five records. The nominal stage succeeds, with terminal test violation $0.51425$. The worst-case, scenario, unbiased KDE, and local shifted stages all report SLSQP failure; their stored controls are useful for diagnosing continuation, but they are not feasible certified solutions. The robust-stage minimum chance margins are respectively approximately $-0.738$, $-0.280$, $-0.0209$, and $-0.0867$ in the solver diagnostics. A terminal test frequency below $0.1$ in a failed stage cannot repair a negative optimization margin, and no pathwise held-out validation was performed. The local experiment therefore supports a statement about solver difficulty and residual statistics, not a claim of successful robust lunar landing.

This distinction also clarifies what a warm start can and cannot do. Continuation reduces abrupt changes in the NLP, but it cannot create feasibility when the transcription, fixed horizon, and risk margins are incompatible. The warm-start paper of \citet{keil2021warm} couples kernel changes to increasing sample sizes and mesh refinement. The local sequence changes the residual estimator while holding a coarse mesh and a fixed sample regime. It is consequently a partial implementation of the continuation idea, not a reproduction of its convergence behavior.

\subsection{Scope of the formulation}

The formulation is adequate for a controlled comparison of residual estimators in low dimension. It is not adequate for a claim about continuous-time path safety, because nodewise constraints do not bound the inter-node maximum. It is not adequate for a physical fuel comparison until the quadrature is made consistent with the defect control parameterization. Finally, the finite-dimensional KDE constraint inherits the statistical qualifications developed earlier: convergence of a KDE sequence under assumptions, as in \citet{schuster2025}, does not give a finite-sample confidence statement for the particular 250-sample local run. These limits determine the evidence hierarchy used in Chapters~8 and~9.
